\documentclass[11pt]{article}

\usepackage[a4paper,margin=2.5cm]{geometry}
\usepackage[ngerman]{babel}
\usepackage{amsmath,amssymb,amsthm,amsfonts}
\usepackage{mathtools}
\usepackage{hyperref}
\usepackage{enumitem}

\title{Priml\"ucken als Rotationen auf einem zyklischen Vier-Zustands-Raum}
\author{Thomas Hoffbauer}
\date{}

\newtheorem{theorem}{Satz}
\newtheorem{lemma}{Lemma}
\newtheorem{corollary}{Korollar}
\newtheorem{proposition}{Proposition}

\theoremstyle{definition}
\newtheorem{definition}{Definition}
\newtheorem{remark}{Bemerkung}

\begin{document}

\maketitle

\begin{abstract}
Wir betrachten Priml\"ucken $g_n=p_{n+1}-p_n$ zwischen aufeinanderfolgenden Primzahlen $>3$ modulo $12$. Die zul\"assigen Restklassen sind $\{1,5,7,11\}=(\mathbb{Z}/12\mathbb{Z})^\times$ mit $(\mathbb{Z}/12\mathbb{Z})^\times\cong C_2\times C_2$; verwendet wird jedoch nur eine zus\"atzlich gew\"ahlte zyklische Ordnung dieser vier Klassen. Der Hauptsatz zeigt, dass das zugeh\"orige Rotationsinkrement auf einem gerichteten Kreisgraphen $C_4$ ausschlie{\ss}lich von $g_n\bmod 12$ abh\"angt. Daraus folgt eine wohldefinierte Abbildung der sechs geraden L\"uckenklassen modulo $12$ auf vier Rotationszust\"ande. Die Konstruktion liefert damit eine exakte kombinatorische Kodierung von Priml\"ucken als diskrete Rotationsdynamik.
\end{abstract}

\section{Einleitung}

Priml\"ucken werden traditionell als arithmetische Gr\"o{\ss}en untersucht,
\[
g_n = p_{n+1}-p_n,
\]
wobei $p_n$ und $p_{n+1}$ aufeinanderfolgende Primzahlen sind.

Das Ziel dieser Notiz ist zu zeigen, dass dieselbe Information eine nat\"urliche geometrische Interpretation besitzt.
Anstatt Priml\"ucken als Abst\"ande auf der Zahlengeraden zu betrachten, interpretieren wir sie als Erzeuger von Rotationen auf einem zyklischen Vier-Zustands-System, das aus den zul\"assigen Restklassen modulo zw\"olf gewonnen wird.

Der mathematische Anspruch dieser Darstellung liegt in einer sauberen, exakten Umcodierung, nicht in einer neuen tiefen Struktur der Primzahlen. Insbesondere wird keine neue Gruppenstruktur eingef\"uhrt; vielmehr wird auf einer endlichen Vierermenge eine zus\"atzliche zyklische Ordnung gew\"ahlt, die die L\"uckenstatistik in eine dynamische Sprache \"ubersetzt.

\section{Der Restklassen-Zustandsraum}

Sei
\[
U_{12}
=
(\mathbb{Z}/12\mathbb{Z})^\times
=
\{1,5,7,11\}.
\]
Jede Primzahl $p>3$ geh\"ort genau zu einer dieser Restklassen.
Als Gruppe gilt $U_{12}\cong C_2\times C_2$, also ist $U_{12}$ nicht zyklisch.
Die Verwendung eines Zyklus bezieht sich hier nur auf eine zus\"atzliche Ordnung dieser vier Elemente.

\begin{definition}
Definiere die Zustandsabbildung
\[
\rho:U_{12}\rightarrow \mathbb{Z}_4
\]
durch
\[
\rho(1)=0,\qquad
\rho(5)=1,\qquad
\rho(7)=2,\qquad
\rho(11)=3.
\]
\end{definition}

Damit bestimmt jede Primzahl $p>3$ einen Zustand
\[
s(p)=\rho(p\bmod 12)\in\mathbb{Z}_4.
\]
Geometrisch identifiziert dies die zul\"assigen Restklassen mit dem gerichteten Kreisgraphen
\[
C_4:\qquad 0\rightarrow1\rightarrow2\rightarrow3\rightarrow0.
\]
\"Aquivalent dazu w\"ahlen wir auf den urspr\"unglichen Restklassen die zyklische Ordnung
\[
1\to5\to7\to11\to1.
\]

Aus historischen Gr\"unden verwenden wir gelegentlich auch die Bezeichnungen
\[
0=E,\qquad 1=A,\qquad 2=B,\qquad 3=C.
\]

\section{Priml\"ucken und Rotationsinkremente}

Seien $p_n<p_{n+1}$ aufeinanderfolgende Primzahlen und
\[
g_n=p_{n+1}-p_n.
\]
Dann gilt f\"ur $p_n>3$:
\[
g_n\equiv 0,2,4,6,8,10 \pmod{12}.
\]

\begin{definition}
Das zu $g_n$ geh\"orige Rotationsinkrement ist
\[
R_n=s(p_{n+1})-s(p_n)\pmod 4.
\]
\end{definition}

\begin{lemma}[Translationen auf den zul\"assigen Restklassen]
Sei $c:\mathbb{Z}_4\to U_{12}$ durch
\[
c(0)=1,\qquad c(1)=5,\qquad c(2)=7,\qquad c(3)=11
\]
gegeben (also $\rho=c^{-1}$).
F\"ur $d\in\{0,2,4,6,8,10\}$ und
\[
T_d(x)\equiv x+d\pmod{12}
\]
ist $T_d$ eine Permutation von $U_{12}$.
Insbesondere ist
\[
\sigma_d:=\rho\circ T_d\circ c:\mathbb{Z}_4\to\mathbb{Z}_4
\]
eine Permutation von $\mathbb{Z}_4$.
\end{lemma}

\begin{proof}
Da $d$ gerade ist, bleibt $U_{12}=\{1,5,7,11\}$ unter $T_d$ invariant.
Auf der endlichen Menge $U_{12}$ folgt daraus sofort die Bijektivit\"at von $T_d$.
Da $c$ und $\rho$ Bijektionen sind, ist auch
$\sigma_d=\rho\circ T_d\circ c$ bijektiv, also eine Permutation von $\mathbb{Z}_4$.
\end{proof}

\begin{proposition}
Sei $\rho:U_{12}\to\mathbb{Z}_4$ die gew\"ahlte zyklische Nummerierung und
f\"ur $d\in\{0,2,4,6,8,10\}$ die Abbildung
\[
T_d:U_{12}\to U_{12},\qquad T_d(x)\equiv x+d\pmod{12}.
\]
Dann ist jedes $T_d$ eine Permutation von $U_{12}$, und es gibt genau ein $r(d)\in\mathbb{Z}_4$ mit
\[
\rho(T_d(x))\equiv \rho(x)+r(d)\pmod 4
\quad\text{f\"ur alle }x\in U_{12}.
\]
Explizit gilt
\[
r(0)=0,\quad r(2)=r(4)=1,\quad r(6)=2,\quad r(8)=r(10)=3.
\]
\end{proposition}

\begin{proof}
Nach dem vorigen Lemma ist $T_d$ eine Permutation von $U_{12}$.
Die Werte von $r(d)$ ergeben sich durch direkte Auswertung von $T_d$ auf den vier Klassen:
\[
\begin{array}{c|c}
d \bmod 12 & \text{\"Ubergangsmuster } s\mapsto s+r(d)\ (\bmod 4)\\
\hline
0   & 0\to0,\ 1\to1,\ 2\to2,\ 3\to3\\
2   & 0\to1,\ 1\to2,\ 2\to3,\ 3\to0\\
4   & 0\to1,\ 1\to2,\ 2\to3,\ 3\to0\\
6   & 0\to2,\ 1\to3,\ 2\to0,\ 3\to1\\
8   & 0\to3,\ 1\to0,\ 2\to1,\ 3\to2\\
10  & 0\to3,\ 1\to0,\ 2\to1,\ 3\to2
\end{array}
\]
Damit ist die Formel f\"ur $r(d)$ gezeigt.
\end{proof}

\section{Darstellungssatz f\"ur Rotationen}

\begin{theorem}
Seien $p_n<p_{n+1}$ aufeinanderfolgende Primzahlen gr\"o{\ss}er als drei.
Dann h\"angt das Rotationsinkrement
\[
R_n=s(p_{n+1})-s(p_n)\pmod 4
\]
nur von der Restklasse der Priml\"ucke modulo zw\"olf ab.
Genauer gilt
\[
R_n=\Phi(g_n\bmod 12),
\]
wobei
\[
\Phi:\{0,2,4,6,8,10\}\rightarrow\mathbb{Z}_4
\]
durch
\[
\Phi(0)=0,\qquad
\Phi(2)=1,\qquad
\Phi(4)=1,\qquad
\Phi(6)=2,\qquad
\Phi(8)=3,\qquad
\Phi(10)=3.
\]
\end{theorem}

\begin{proof}
Setze $d=g_n\bmod 12\in\{0,2,4,6,8,10\}$.
Nach der vorigen Proposition ist die Wirkung von $x\mapsto x+d$ auf $U_{12}$
eine Permutation mit eindeutigem Inkrement $r(d)\in\mathbb{Z}_4$, so dass
\[
s(p_{n+1})-s(p_n)\equiv r(d)\pmod 4.
\]
Damit ist $R_n$ bereits eine Funktion nur von $d=g_n\bmod 12$, n\"amlich $R_n=r(d)$.
Mit der expliziten Tabelle aus der Proposition folgt:
\[
g_n\equiv0 \Rightarrow R_n=0,\qquad
g_n\equiv2,4 \Rightarrow R_n=1,\qquad
g_n\equiv6 \Rightarrow R_n=2,\qquad
g_n\equiv8,10 \Rightarrow R_n=3.
\]
Also h\"angt $R_n$ nur von $g_n \bmod 12$ ab.
\end{proof}

\begin{remark}[Informations\"aquivalenz]
Der Hauptsatz ist eine exakte Umcodierung:
Die Gr\"o{\ss}e $R_n$ enth\"alt keine zus\"atzliche arithmetische Information
gegen\"uber $g_n\bmod 12$, sondern kodiert dieselbe Information als
orientierten Schritt auf $C_4$.
\end{remark}

\section{Interpretation}

\begin{remark}
Der Begriff \emph{Rotation} wird hier graphentheoretisch verwendet:
Gemeint ist die orientierte Verschiebung auf dem gerichteten Kreisgraphen $C_4$,
nicht eine geometrische Rotation im euklidischen Raum.
\end{remark}

\begin{proposition}[Quaternionische Darstellung mit Hamiltonzahlen]
Sei $\mathbb{H}$ die Algebra der Hamiltonzahlen und
\[
Q_4:=\{1,\mathbf{i},-1,-\mathbf{i}\}\subset\mathbb{H}.
\]
Dann ist $Q_4$ eine zyklische Gruppe der Ordnung $4$ bez\"uglich der Multiplikation, erzeugt von $\mathbf{i}$.
Die Abbildung
\[
\eta:\mathbb{Z}_4\to Q_4,\qquad \eta(r):=\mathbf{i}^r
\]
ist ein Gruppenisomorphismus.
\end{proposition}

\begin{proof}
Aus $\mathbf{i}^2=-1$ und $\mathbf{i}^4=1$ folgt
\[
Q_4=\{1,\mathbf{i},\mathbf{i}^2,\mathbf{i}^3\}=\{1,\mathbf{i},-1,-\mathbf{i}\}.
\]
Damit ist $Q_4$ zyklisch von Ordnung $4$. Ferner gilt f\"ur $a,b\in\mathbb{Z}_4$:
\[
\eta(a+b)=\mathbf{i}^{a+b}=\mathbf{i}^a\mathbf{i}^b=\eta(a)\eta(b),
\]
also ist $\eta$ ein Homomorphismus. Bijektivit\"at ist wegen gleicher endlicher
Kardinalit\"at unmittelbar.
\end{proof}

Der Satz liefert eine exakte Faktorisierung
\[
\text{Priml\"ucken}
\longrightarrow
\mathbb{Z}/12\mathbb{Z}
\longrightarrow
C_4.
\]
Damit erzeugt jede Priml\"ucke eine eindeutige Rotation auf dem Kreisgraphen.
\"Uber $\eta$ kann man dieselbe Dynamik \"aquivalent in Hamiltonzahlen schreiben:
\[
q_n:=\eta(s(p_n))\in Q_4,\qquad
q_{n+1}=q_n\cdot \mathbf{i}^{R_n}.
\]
Die Rotationsinkremente erscheinen damit als multiplikative Faktoren in der
zyklischen Quaternion-Untergruppe $Q_4$.
Die nichtkommutative Gesamtstruktur von $\mathbb{H}$ wird dabei nicht benutzt.

Die Zuordnung lautet:
\[
\begin{array}{c|c|c}
g_n\bmod 12 & R_n & \text{Interpretation}\\
\hline
0 & 0 & \text{Keine Bewegung}\\
2,4 & 1 & \text{Vorw\"artsrotation}\\
6 & 2 & \text{Halbdrehung}\\
8,10 & 3 & \text{R\"uckw\"artsrotation}
\end{array}
\]

Folglich induziert die Primzahlfolge einen diskreten Rotationsfluss
\[
(s(p_n))_{n\ge 1}
\]
auf dem Kreisgraphen $C_4$.

\section{Rotationsw\"orter f\"ur Primkonstellationen}

\begin{definition}
Sei
\[
\mathcal{H}=(h_0,h_1,\dots,h_{k-1}),\qquad h_0<h_1<\dots<h_{k-1},
\]
ein Tupel ganzer Zahlen. Falls f\"ur ein $p>3$ alle Zahlen
\[
p+h_0,\;p+h_1,\;\dots,\;p+h_{k-1}
\]
prim sind, nennen wir
\[
(p+h_0,\dots,p+h_{k-1})
\]
eine \emph{realisierte Primkonstellation} vom Typ $\mathcal{H}$.
Ihr \emph{Rotationswort} ist
\[
W_{\mathcal{H}}:=(R_0,\dots,R_{k-2})\in\mathbb{Z}_4^{\,k-1},
\]
wobei
\[
R_i=s(p+h_{i+1})-s(p+h_i)\pmod 4.
\]
\end{definition}

\begin{proposition}
F\"ur eine realisierte Primkonstellation vom Typ
\[
\mathcal{H}=(h_0,\dots,h_{k-1})
\]
h\"angt das Rotationswort nur von den Differenzen
\[
\Delta_i:=h_{i+1}-h_i\qquad (0\le i\le k-2)
\]
modulo $12$ ab. Genauer gilt
\[
R_i=\Phi(\Delta_i\bmod 12)\qquad (0\le i\le k-2).
\]
Insbesondere ist das Rotationswort innerhalb eines festen Konstellationstyps
f\"ur jede Realisierung identisch.
\end{proposition}

\begin{proof}
F\"ur jedes $i$ haben wir
\[
p+h_{i+1}=(p+h_i)+\Delta_i.
\]
Reduziert modulo $12$ ist dies genau die Translation
\[
x\longmapsto x+\Delta_i \pmod{12}
\]
auf $U_{12}=\{1,5,7,11\}$. Nach der Proposition zu den Translationen $T_d$
h\"angt das zugeh\"orige Inkrement
\[
R_i=s(p+h_{i+1})-s(p+h_i)\pmod 4
\]
nur von $\Delta_i\bmod 12$ ab und ist gleich $\Phi(\Delta_i\bmod 12)$.
Dies gilt f\"ur alle $i$, also folgt die Behauptung.
\end{proof}

\begin{corollary}
F\"ur den Primzahlvierlingstyp
\[
\mathcal{H}=(0,2,6,8)
\]
ist
\[
W_{\mathcal{H}}=(1,1,1).
\]
\end{corollary}

\begin{proof}
Hier sind die Differenzen
\[
(\Delta_0,\Delta_1,\Delta_2)=(2,4,2).
\]
Nach der Definition von $\Phi$ gilt
\[
\Phi(2)=\Phi(4)=1,
\]
also
\[
W_{\mathcal{H}}=(\Phi(2),\Phi(4),\Phi(2))=(1,1,1).
\]
\end{proof}

\section{Dynamische Sicht}

Die Primzahlfolge
\[
(p_n)
\]
bestimmt daher eine Trajektorie
\[
(s(p_n))
\]
auf dem endlichen Zustandsraum $C_4$.
Priml\"ucken wirken als Bewegungserzeuger, w\"ahrend Primzahlkonstellationen als koh\"arente Rotationsmuster erscheinen.

Dies ergibt eine geometrische Neuformulierung der Priml\"uckenstatistik:
\[
\boxed{
\text{Priml\"ucken modulo }12
\Longleftrightarrow
\text{Rotationen auf }C_4.
}
\]

\subsection*{Rotationsbeschreibung des ehemaligen Miller-Effekts}

Der \emph{ehemalige Miller-Effekt} l\"asst sich in der vorliegenden Sprache als
Orientierungsdrift des Rotationsflusses formulieren:
Vorw\"artsinkremente $R_n=1$ (modulo $4$) treten h\"aufiger auf als
R\"uckw\"artsinkremente $R_n=3$.

Eine \"aquivalente Beschreibung liefern die Gr\"o{\ss}en
\[
M_N:=\frac{\#\{1\le n\le N:\,R_n=1\}}{\#\{1\le n\le N:\,R_n=3\}}
\]
und
\[
E_N:=\#\{1\le n\le N:\,R_n=1\}-\#\{1\le n\le N:\,R_n=3\}.
\]
Dann entspricht der Effekt den Aussagen
\[
M_N>1
\quad\Longleftrightarrow\quad
E_N>0
\]
auf gro{\ss}en endlichen Bereichen.

Damit ist der ehemalige Miller-Effekt nicht als neue Restklasseninformation zu
verstehen, sondern als statistische Asymmetrie in der symbolisch-dynamischen
Auswertung der Rotationsfolge $(R_n)$.

\section{Ausblick}

Ein naheliegender n\"achster Schritt ist eine nichttriviale dynamische Aussage \"uber die induzierte Folge $(R_n)$. Als erste Leitgr\"o{\ss}e bietet sich das Verh\"altnis
\[
M_N
:=
\frac{\#\{1\le n\le N:\,R_n=1\}}
{\#\{1\le n\le N:\,R_n=3\}}.
\]
\emph{Offene Alternative.} Entweder gilt
\[
\lim_{N\to\infty} M_N = 1,
\]
oder es existiert ein Grenzwert $L\neq1$ mit
\[
\lim_{N\to\infty} M_N = L.
\]
Als empirisch vorsichtigere Zwischenhypothese kann man
\[
M_N>1
\]
auf gro{\ss}en endlichen Bereichen testen.

Methodisch zentral ist dabei:
Nicht die Einzelgr\"o{\ss}e $R_n$, sondern endliche Rotationsw\"orter
\[
(R_n,R_{n+1},\dots,R_{n+\ell})
\]
sind der vielversprechende Tr\"ager m\"oglicher neuer Struktur.
F\"ur $k\ge 1$ sei dazu
\[
\mathcal{L}_k
:=
\left\{
(R_n,R_{n+1},\dots,R_{n+k-1})\in\mathbb{Z}_4^k:\ n\ge 1
\right\}
\]
die Menge der beobachteten Rotationsw\"orter der L\"ange $k$.
Die Struktur von $\mathcal{L}_k$ (zul\"assige und verbotene W\"orter)
ist eine zentrale Forschungsfrage.

Dar\"uber hinaus ergeben sich die folgenden konkreten Forschungsvorschl\"age:
\begin{enumerate}[label=\textbf{F\arabic*:},leftmargin=*,itemsep=0.35em]
\item \textbf{\"Ubergangsmatrix und Markov-N\"aherung.}
Definiere f\"ur $i,j\in\mathbb{Z}_4$ die empirischen \"Ubergangszahlen
\[
N_{ij}(X):=\#\{n:\ p_n\le X,\ s(p_n)=i,\ s(p_{n+1})=j\},
\]
und die normierte Matrix $P_X=(p_{ij}(X))$ mit
\[
p_{ij}(X):=\frac{N_{ij}(X)}{\sum_{k}N_{ik}(X)}.
\]
Zu untersuchen ist die Konvergenz von $P_X$ und eine m\"ogliche Asymmetrie zwischen Vorw\"arts- und R\"uckw\"arts\"uberg\"angen.

\item \textbf{H\"aufigkeiten der Rotationsklassen.}
Setze
\[
\nu_r(X):=\#\{n:\ p_n\le X,\ R_n=r\}\qquad (r\in\mathbb{Z}_4).
\]
Die zentralen Fragen sind Grenzverh\"altnisse wie $\nu_1(X)/\nu_3(X)$, sowie Fehlerterme gegen ein Referenzmodell ohne Bias.

\item \textbf{Vergleich mit Zufallsmodellen.}
Vergleiche die beobachteten Werte von $M_N$, $\nu_r(X)$ und $P_X$ mit
geeigneten Cram\'er-artigen oder entkoppelten Restklassenmodellen.
Damit l\"asst sich trennen, welche Effekte rein lokal-modular sind und welche auf tieferer Korrelation beruhen.

\item \textbf{Rotationsw\"orter und Konstellationsklassen.}
Verallgemeinere das Vierlings-Korollar auf l\"angere Muster und untersuche:
Welche Rotationsw\"orter sind \"uberhaupt zul\"assig?
Welche W\"orter entsprechen admissiblen $k$-Tupeln?
Welche W\"orter zeigen systematische H\"aufigkeitsabweichungen?
Dies ist der symbolisch-dynamische Kern der Methode.

\item \textbf{Verallgemeinerung des Moduls.}
Ersetze $12$ durch ein allgemeineres Modul $m$ mit fixierten zul\"assigen Restklassen.
Die Frage ist, wann eine analoge Kodierung auf einem endlichen gerichteten Graphen
entsteht und welche dynamischen Gr\"o{\ss}en dann stabil bleiben.
\end{enumerate}

Diese Agenda macht aus der vorliegenden Kodierung ein systematisches Forschungsprogramm,
in dem numerische Evidenz, kombinatorische Struktur und analytische Heuristik unmittelbar zusammenwirken.

\section{Schluss}

Wir haben gezeigt, dass die zul\"assigen Restklassen der Primzahlen modulo zw\"olf eine nat\"urliche Darstellung als Vier-Zustands-Raum mit gew\"ahlter zyklischer Ordnung besitzen.
Unter dieser Darstellung induziert jede Priml\"ucke eine diskrete Rotation, deren Wert nur von der Restklasse der L\"ucke modulo zw\"olf abh\"angt.
Die resultierende Entsprechung ist eine exakte kombinatorische Umcodierung von Priml\"uckenklassen in ein diskretes Rotationssystem auf einem endlichen Kreisgraphen.

\end{document}